%% ma: L12-B14-0139
%% lop: 12
%% chuong: 5
%% bai: 14
%% dang: 12.14.12
%% loai: DS
%% muc: [NB, NB, TH, VD]
%% dap_an: "SĐĐĐ"
%% nguon: "Ảnh giáo viên gửi 25/11/2026 (tài liệu ôn tập Phương trình mặt phẳng), Phần II Câu 4 (điểm A(2;1;1), B(-1;-2;-3), mặt phẳng (Q): x+y+z=0)"
%% kien_thuc: "Bài 7 (tọa độ vectơ), Bài 14 (tích có hướng, vectơ pháp tuyến, phương trình mặt phẳng)"
%% trang_thai: chinh-thuc
%% ghi_chu: "Đáp án tự giải (sympy). Ý a) sai vì {AB}=(-3;-3;-4)."
\begin{ex}
Trong không gian $Oxyz$, cho hai điểm $A(2;1;1)$, $B(-1;-2;-3)$ và mặt phẳng $(Q)\colon x+y+z=0$.
\choiceTF
{$\overrightarrow{AB}=(-3;-3;4)$.}
{\True Một vectơ pháp tuyến của $(Q)$ là $\vec n_Q=(1;1;1)$.}
{\True $\left[\overrightarrow{AB},\vec n_Q\right]=(1;-1;0)$.}
{\True Mặt phẳng qua $A,B$ và vuông góc với $(Q)$ có phương trình là $x-y-1=0$.}
\loigiai{a) $\overrightarrow{AB}=(-3;-3;-4)$. Sai.
b) Đúng.
c) $\left[\overrightarrow{AB},\vec n_Q\right]=\big((-3)\cdot1-(-4)\cdot1;\ (-4)\cdot1-(-3)\cdot1;\ (-3)\cdot1-(-3)\cdot1\big)=(1;-1;0)$. Đúng.
d) Mặt phẳng qua $A,B$ và vuông góc $(Q)$ nhận $\left[\overrightarrow{AB},\vec n_Q\right]=(1;-1;0)$ làm vectơ pháp tuyến: $(x-2)-(y-1)=0\Leftrightarrow x-y-1=0$ (điểm $B$ cũng thuộc mặt phẳng). Đúng.}
\end{ex}
